%%--------------------------------------------------------------------
%1--------------------------------------------------------------------
\vspace*{\fill}
\begin{glyph}{Beautiful (美)}{beautiful}\label{glyph:beautiful}
All things beautiful.
\end{glyph}

\begin{comps}
\comprtwocone & 羊 + 大 \\\hline
    \comprthreecone & Sheep (\ref{glyph:sheep}) + Large (\ref{glyph:large}) \\\hline
\end{comps}

\begin{remglyph}
    \hooglig{Beautiful} \remcom{large} \remcom{sheep}.\\\hline
\end{remglyph}

\begin{prons}
\pronsrtwocone & \textbf{ビ (BI)} & \textbf{うつく (utsuku)}\\\hline
\pronsrthreecone & --- & --- \\\hline
\end{prons}

\begin{remprons}
The \comon{big} \hooglig{beautiful} \comkun{boots coupon}.\\\hline
\end{remprons}

%{\middel}
% &  ()\\\hline
\begin{somme}
美しい & \textit{beautiful} & うつく{\middel}しい (utsuku{\middel}shii)\\\hline
美人 & beautiful + person = \textit{beautiful woman} &　ビ{\middel}ジン (BI{\middel}JIN) \\\hline
美化 & beautiful + change = \textit{beautification} & ビ{\middel}カ (BI{\middel}KA)\\\hline
    \notyet{美学} & beautiful + study = \textit{aesthetics} & ビ{\middel}ガク (BI{\middel}GAKU)\\\hline
\end{somme}

\begin{irrr}
    \notyet{美味しい} & beautiful + taste = \textit{delicious} & おい{\middel}しい (oi{\middel}shii)\\\hline	
\end{irrr}

\begin{sentence}
\ruby{\underline{山}}{やま}の\ruby{\underline{木}}{き}が\ruby{\underline{美しい}}{うつくしい}です。\\\hline
yama no ki ga utsuku{\middel}shii desu.\\\hline
(mountain) (trees) (beautiful)\\\hline
=\textit{The mountain trees are beautiful.}\\\hline
\end{sentence}
\end{multicols}
\vspace*{-\baselineskip}
\vhrulefill{5pt}
%%--------------------------------------------------------------------
%2--------------------------------------------------------------------
\begin{multicols}{2}
\begin{glyph}{Island (島)}{island}\label{glyph:island}
Island.
\end{glyph}

\begin{comps}
\comprtwocone & 鳥 + 山  \\\hline
    \comprthreecone &Bird (\ref{glyph:bird}) + Mountain (\ref{glyph:mountain})  \\\hline
\end{comps}

\begin{remglyph}
    This \remcom{bird} lives only on the \remcom{mountains} of a particular
    \hooglig{island}.\\\hline
\end{remglyph}

\begin{prons}
\pronsrtwocone & \textbf{トウ (T\={O})} & \textbf{しま (shima)}\\\hline
\pronsrthreecone & --- & --- \\\hline
\end{prons}

\begin{remprons}
\comkun{She mugs} \comon{tortured} \hooglig{islanders}.\\\hline
\end{remprons}

%{\middel}
% &  ()\\\hline
\begin{somme}
    島 & \textit{island} &しま (shima) \\\hline
    小島 & small + island = \textit{small island} & こ{\middel}じま (ko{\middel}jima)\\\hline
    島国 & island + country = \textit{island nation} & しま{\middel}ぐに (shima{\middel}guni)\\\hline
    半島 & half + island = \textit{Peninsula} & ハン{\middel}トウ (HAN{\middel}T\={O})\\\hline
    本島 & main + island = \textit{Main island} & ホン{\middel}トウ (HON{\middel}T\={O})\\\hline
\end{somme}
\begin{sentence}
    \underline{オウストラリア}と\ruby{\underline{日本}}{ニホン}\ruby{は}{わ}\ruby{\underline{島国}}{しまぐに}です。\\\hline
    \={O}sutoraria to NI{\middel}HON wa shima{\middel}guni desu.\\\hline
    (Australia) (Japan) (island nation) \\\hline
    =\textit{Australia and Japan are island nations.}\\\hline
\end{sentence}
\end{multicols}
\vspace*{-\baselineskip}
\vhrulefill{5pt}
\vspace*{\fill}
\pagebreak
%%--------------------------------------------------------------------
%3--------------------------------------------------------------------
\vspace*{\fill}
\begin{multicols}{2}
\begin{glyph}{Tiny (寸)}{tiny}\label{glyph:tiny}
Tiny.
\end{glyph}

\begin{comps}
\comprtwocone & 寸 \\\hline
    \comprthreecone & small/inch (\ref{glyph:tiny}) \\\hline
\end{comps}

\begin{remglyph}
    An \remcom{inch} is \hooglig{small}.\\\hline
\end{remglyph}

\begin{prons}
\pronsrtwocone & \textbf{スン (SUN)} & ---\\\hline
\pronsrthreecone & --- & --- \\\hline
\end{prons}

\begin{remprons}
    \hooglig{Tiny} \comon{soon}.\\\hline
\end{remprons}

%{\middel}
% &  ()\\\hline
\begin{somme}
    一寸  & one + tiny = \textit{a tiny bit}  & イッ{\middel}スン (IS{\middel}SUN) \\\hline
    \notyet{寸前}  & tiny + before = \textit{immediately before}  & スン{\middel}ゼン
    (SUN{\middel}ZEN) \\\hline
\end{somme}

\begin{sentence}
    \ruby{\underline{出}}{で}\underline{かける}{\;}\ruby{\underline{寸前}}{スンゼン}に\ruby{\underline{高木}}{たかぎ}\underline{さん}\ruby{を}{お}\ruby{\underline{見}}{み}\underline{ました}。\\\hline
    de{\middel}kakeru SUN{\middel}ZEN ni Taka{\middel}gi-san o mi{\middel}mashita.\\\hline
    (leave) (immediately before) (Takagi-san) (saw)\\\hline
    =\textit{I saw Takagi-san just before I left.} \\\hline
\end{sentence}
\end{multicols}
\vspace*{-\baselineskip}
\vhrulefill{5pt}
%%--------------------------------------------------------------------
%4--------------------------------------------------------------------
\begin{multicols}{2}
\begin{glyph}{Temple (寺)}{temple}\label{glyph:temple}
    (Buddhist) Temple.
\end{glyph}

\begin{comps}
\comprtwocone &  寸 + 土 \\\hline
    \comprthreecone & Tiny (\ref{glyph:tiny}) + Earth (\ref{glyph:earth}) \\\hline
\end{comps}

\begin{remglyph}
    A \hooglig{temple} is \remcom{tiny} compared to the entire \remcom{earth}.\\\hline
\end{remglyph}

\begin{prons}
\pronsrtwocone & \textbf{ジ (JI)} & \textbf{てら (tera)}\\\hline
\pronsrthreecone & ---  & ---  \\\hline
\end{prons}

\begin{remprons}
    The \comon{jig} is up.  You will not \comkun{terraform} this
    \hooglig{temple}.\\\hline
\end{remprons}

%{\middel}
% &  ()\\\hline
\begin{somme}
    寺  & \textit{temple}  & てら (tera) \\\hline
    寺田さん  & temple + main = \textit{Terada-san}  & てら{\middel}だ{\middel}さん (tera{\middel}da-san) \\\hline
    山寺  & mountain + temple = \textit{mountain temple}  & やま{\middel}でら (yama{\middel}dera) \\\hline
    古寺  & old + temple = \textit{old temple}  & ふる{\middel}でら (furu{\middel}dera) コ{\middel}ジ (KO{\middel}JI) \\\hline
\end{somme}

\begin{sentence}
    \ruby{\underline{山寺}}{やまでら}の\ruby{\underline{古}}{ふる}\underline{い}{\;}\ruby{\underline{門}}{モン}の\ruby{\underline{下}}{した}で\ruby{\underline{休}}{やす}\underline{みました}。\\\hline
    yama{\middel}dera no furu{\middel}i MON no shita de yasu{\middel}mimashita.\\\hline
    (mountain temple) (old) (gate) (under) (rested) \\\hline
    =\textit{(We) rested under the old gate of the mountain temple.}\\\hline
\end{sentence}
\end{multicols}
\vspace*{-\baselineskip}
\vhrulefill{5pt}
\vspace*{\fill}
\pagebreak
%%--------------------------------------------------------------------
%5--------------------------------------------------------------------
\begin{multicols}{2}
\begin{glyph}{Precede (先)}{precede}\label{glyph:precede}
All meanings derived from the general sense of ``preceding'', or the idea of an
    object ``ahead'' in a literal or figurative way.\\\hline
    Depending on the context in which it appears, this kanji can signify the
    ``tip'' or ``point'' of something, the general future (or the past, when used as
    in the third compound), or a destination.\\\hline
    It can also refer to people in the sense of ``the other party''.
\end{glyph}

\begin{comps}
\comprtwocone &  牛 + 儿\\\hline
    \comprthreecone & Cow (\ref{glyph:cow}) + Person (\ref{glyph:person}) \\\hline
\end{comps}

\begin{remglyph}
    The \remcom{cow} \hooglig{preceded} the \remcom{person}.\\\hline
\end{remglyph}

\begin{prons}
\pronsrtwocone & \textbf{セン (SEN)} & \textbf{さき (saki)}\\\hline
\pronsrthreecone & --- & --- \\\hline
\end{prons}

\begin{remprons}
    It's \comkun{sucky} that the \hooglig{preceding} rules are so
    \comon{central} to this school.\\\hline
\end{remprons}

%{\middel}
% &  ()\\\hline
\begin{somme}
    先  & \textit{precede}  & さき (saki) \\\hline
    先生  & precede + life = \textit{teacher}  & セン{\middel}セイ (SEN{\middel}SEI) \\\hline
    先月  & precede + moon (month) = \textit{last month}  & セン{\middel}ゲツ (SEN{\middel}GETSU) \\\hline
    口先  & mouth + precede = \textit{lip service}  & くち{\middel}さき (kuchi{\middel}saki)\\\hline
\end{somme}

\begin{sentence}
    \ruby{\underline{先生}}{センセイ}\ruby{は}{わ}\ruby{\underline{冬休}}{ふゆやす}\underline{み}に\ruby{\underline{日本}}{ニホン}{\;}\underline{アルプス}\ruby{へ}{え}\ruby{\underline{行}}{い}\underline{きました}。\\\hline
    SEN{\middel}SEI wa fuyu yasu{\middel}mi ni NI{\middel}HON arupusu e i{\middel}kimashita.\\\hline
    (teacher) (winter vacation) (Japanese Alps) (went)\\\hline
    =\textit{(Our) teacher went to the Japanese Alps for winter
    vacation.}\\\hline
\end{sentence}
\end{multicols}
%%--------------------------------------------------------------------
%6--------------------------------------------------------------------
\begin{multicols}{2}
\begin{glyph}{Open (開)}{open}\label{glyph:open}
    Open.\\\hline
    Secondary meanings include ``starting'' and the idea of ``developing''
    things (in the sense of opening or reclaiming land for example.)
\end{glyph}

\begin{comps}
\comprtwocone & 門 + 廾 + 一 \\\hline
    \comprthreecone & Gate (\ref{glyph:gate}) + two hands + One (\ref{glyph:one}) \\\hline
\end{comps}

\begin{remglyph}
    It takes \remcom{two hands} to \hooglig{open} the \remcom{gate} in
    \remcom{one} go.\\\hline
\end{remglyph}

\begin{prons}
\pronsrtwocone & \textbf{カイ (KAI)} & \textbf{あ; ひら (a; hira)}\\\hline
\pronsrthreecone & --- & --- \\\hline
    \rye{3}{In the general physical of ``to open'', あく and あける function as
    a normal intrasitive/transitive verb pair.\\\hline
    ひらく is more concerned with a metaphorical or non-physical sense of
    ``opening'', and appears primarily in expressions (in the ``opening'', for
    example, of a conference, of new lands, or of cherry blossoms).\\\hline
    Expressing the sense of a shop's front door being open would thus call for
    the verb あく, while conveying the idea of opening a business would require
    ひらく.\\\hline
    In the first position, this kanji invariably takes the onyomi.}
\end{prons}

\begin{remprons}
    \comkun{He rustled} through the \comkun{updates} in hopes of \hooglig{opening} new possibilities to \comon{kite-surfing}.\\\hline
\end{remprons}

%{\middel}
% &  ()\\\hline
\begin{somme}
    開く (intr)  & \textit{to be open}  & あ{\middel}く (a{\middel}ku) \\\hline
    開ける (tr)  & \textit{to open (something)}  & あ{\middel}ける (a{\middel}keru) \\\hline
    開く (intr/tr)  & \textit{to open; to develop}  & ひら{\middel}く (hira{\middel}ku) \\\hline
    開化  & open + change = \textit{becoming civilized}  & カイ{\middel}カ (KAI{\middel}KA) \\\hline
    開国  & open + country = \textit{to open up a country}  & カイ{\middel}コク (KAI{\middel}KOKU) \\\hline
    \notyet{開会}  & open + meet = \textit{to open a meeting}  & カイ{\middel}カイ (KAI{\middel}KAI) \\\hline
    \notyet{公開}  & public + open = \textit{open to the public}  & コ{\middel}ウカイ (K\={O}{\middel}KAI) \\\hline
\end{somme}

\begin{sentence}
    \ruby{\underline{水門}}{スイモン}\ruby{を}{お}\ruby{\underline{開}}{あ}\underline{けて}{\;}\ruby{\underline{下}}{くだ}\underline{さい}。\\\hline
    SUI{\middel}MON o a{\middel}kete kuda{\middel}sai.\\\hline
    (sluice gate) (open) (please)\\\hline
    =\textit{Please open the sluice gate.}\\\hline
\end{sentence}
\end{multicols}
\pagebreak
%%--------------------------------------------------------------------
%7--------------------------------------------------------------------
\vspace*{\fill}
\begin{multicols}{2}
\begin{glyph}{Bird (鳥)}{bird}\label{glyph:bird}
    Bird.
\end{glyph}

\begin{comps}
\comprtwocone & 鳥 \\\hline
    \comprthreecone & Bird (\ref{glyph:bird}) \\\hline
\end{comps}

\begin{remglyph}
\\\hline
\end{remglyph}

\begin{prons}
    \pronsrtwocone & \textbf{チョウ (CH\={O})} & \textbf{とり (tori)}\\\hline
\pronsrthreecone & --- & --- \\\hline
\end{prons}

\begin{remprons}
    \comkun{Toreador} \comon{chores} are for the \hooglig{birds}.\\\hline
\end{remprons}

%{\middel}
% &  ()\\\hline
\begin{somme}
    鳥  & \textit{bird}  &  とり (tori) \\\hline
    白鳥  & white + bird = \textit{swan}  & ハク{\middel}チョウ (HAKU{\middel}CH\={O}) \\\hline
    小鳥  & small + bird = \textit{small bird}  &  こ{\middel}とり (ko{\middel}tori) \\\hline
    \notyet{夜鳥}  & night + bird = \textit{nocturnal bird}  & ヤ{\middel}チョウ (YA{\middel}CH\={O})  \\\hline
    \notyet{愛鳥家}  & love + bird + house = \textit{bird lover}  & アイ{\middel}チョウ{\middel}カ (AI{\middel}CH\={O}{\middel}KA) \\\hline
\end{somme}

\begin{sentence}
    \underline{あの}{\;}\ruby{\underline{青}}{あお}\underline{い}{\;}\ruby{\underline{小鳥}}{ことり}の\ruby{\underline{名前}}{なまえ}が\ruby{\underline{分}}{わ}\underline{かります}か。\\\hline
    ano ao{\middel}i ko{\middel}tori no na{\middel}mae ga wa{\middel}karimasu ka.\\\hline
    (that) (blue) (little bird) (name) (understand)\\\hline
    =\textit{Do you know the name of that little blue bird?}\\\hline
\end{sentence}
\end{multicols}
\vspace*{-\baselineskip}
\vhrulefill{5pt}
%%--------------------------------------------------------------------
%8--------------------------------------------------------------------
\begin{multicols}{2}
\begin{glyph}{Evening (夕)}{evening}\label{glyph:evening}
Evening.
\end{glyph}

\begin{comps}
\comprtwocone & 夕 \\\hline
    \comprthreecone & Evening (\ref{glyph:evening}) \\\hline
\end{comps}

\begin{remglyph}
\\\hline
\end{remglyph}

\begin{prons}
    \pronsrtwocone & --- & \textbf{ゆう (y\={u})}\\\hline
    \pronsrthreecone & セキ (SEKI) & --- \\\hline
\end{prons}

\begin{remprons}
    \comkun{Youse} look \ucomon{sexy} this \hooglig{evening}.\\\hline
\end{remprons}

%{\middel}
% &  ()\\\hline
\begin{somme}
    夕べ  & \textit{evening}  & ゆう{\middel}べ (y\={u}{\middel}be) \\\hline
    夕日  & evening + sun = \textit{setting sun}  & ゆう{\middel}ひ (y\={u}{\middel}hi) \\\hline
    \notyet{夕立}  & evening + stand = \textit{late afternoon (rain) shower}  & ゆう{\middel}だち (y\={u}{\middel}dachi)\\\hline
\end{somme}

\begin{sentence}
    \ruby{\underline{山}}{やま}{\;}\underline{から}{\;}\ruby{\underline{見}}{み}\underline{る}{\;}\ruby{\underline{夕日}}{ゆうひ}が\ruby{\underline{美}}{うつく}\underline{しい}。\\\hline
    yama kara mi{\middel}ru y\={u}{\middel}hi ga utsuku{\middel}shii.\\\hline
    (mountain) (from) (see) (setting sun) (beautiful)\\\hline
    =\textit{The setting sun is beautiful when seen from the mountain.}\\\hline
\end{sentence}
\end{multicols}
\vspace*{-\baselineskip}
\vhrulefill{5pt}
\vspace*{\fill}
\pagebreak
%%--------------------------------------------------------------------
%9--------------------------------------------------------------------
\vspace*{\fill}
\begin{multicols}{2}
\begin{glyph}{Outside (外)}{outside}\label{glyph:outside}
Outside.\\\hline
The idea of ``foreign'' or ``other'' is also present, as the third example makes
clear.\\\hline
The less common readings convey some secondary meanings of removal and
    disconnection.
\end{glyph}

\begin{comps}
\comprtwocone & 夕 + 卜  \\\hline
    \comprthreecone &Evening (\ref{glyph:evening}) + divination  \\\hline
\end{comps}

\begin{remglyph}
    Go \remcom{outside} to do your \remcom{divination} in the \remcom{evening}.\\\hline
\end{remglyph}

\begin{prons}
\pronsrtwocone & \textbf{ガイ (GAI)} & \textbf{そと (soto)}\\\hline
    \pronsrthreecone &ゲ (GE) &ほか; はず (hoka; hazu) \\\hline
\end{prons}

\begin{remprons}
    \hooglig{Outside} is a \ucomkun{hawker} \comon{guy} that \ucomon{gets} all the \ucomkun{huzz} and buzz with his \comkun{sotto} voce remarks.\\\hline
\end{remprons}

%{\middel}
% &  ()\\\hline
\begin{somme}
    外  & \textit{outside}  & そと (soto) \\\hline
    外見  & outside + see = \textit{external appearance}  & ガイ{\middel}ケン (GAI{\middel}KEN) \\\hline
    外国人  & outside + country + person = \textit{foreigner}  & ガイ{\middel}コク{\middel}ジン (GAI{\middel}KOKU{\middel}JIN) \\\hline
    外出  & outside + exit = \textit{going out}  & ガイ{\middel}シュツ (GAI{\middel}SHUTSU) \\\hline
    \notyet{海外}  & sea + outside = \textit{overseas}  & カイ{\middel}ガイ (KAI{\middel}GAI) \\\hline
\end{somme}

\begin{sentence}
    \ruby{\underline{日本}}{ニホン}には\ruby{\underline{外国人}}{ガイコクジン}が\ruby{\underline{多}}{おお}\underline{い}。\\\hline
    NI{\middel}HON ni wa GAI{\middel}KOKU{\middel}JIN ga \={o}{\middel}i.\\\hline
    (Japan) (foreigners) (many)\\\hline
    =\textit{There are many foreigners in Japan.}\\\hline
\end{sentence}
\end{multicols}
\vspace*{-\baselineskip}
\vhrulefill{5pt}
%%--------------------------------------------------------------------
%10-------------------------------------------------------------------
\begin{multicols}{2}
\begin{glyph}{Tongue (舌)}{tongue}\label{glyph:tongue}
Tongue.
\end{glyph}

\begin{comps}
\comprtwocone &千 + 口  \\\hline
    \comprthreecone & Thousand (\ref{glyph:thousand}) + Mouth (\ref{glyph:mouth}) \\\hline
\end{comps}

\begin{remglyph}
    The \hooglig{tongue} in your \remcom{mouth} can lead to \remcom{thousands} of consequences.\\\hline
\end{remglyph}

\begin{prons}
\pronsrtwocone & --- & \textbf{した (shita)}\\\hline
    \pronsrthreecone &ゼツ (ZETSU)  & --- \\\hline
\end{prons}

\begin{remprons}
    \comkun{She tugged} \ucomon{Zetsu's} \hooglig{tongue}.\\\hline
\end{remprons}

%{\middel}
% &  ()\\\hline
\begin{somme}
    舌  & \textit{tongue}  & した (shita) \\\hline
\end{somme}

\begin{sentence}
    \ruby{\underline{内田}}{うちだ}\underline{さん}\ruby{は}{わ}\ruby{\underline{舌}}{した}に\underline{ピアス}\ruby{を}{お}\underline{しました}。\\\hline
    Uchi{\middel}da-san wa shita ni piasu o shimashita.\\\hline
    (Uchida-san) (togue) (piercing) (did)\\\hline
    =\textit{Uchida-san got her tongue pierced.}\\\hline
\end{sentence}
\end{multicols}
\vspace*{-\baselineskip}
\vhrulefill{5pt}
\vspace*{\fill}
\pagebreak
%%--------------------------------------------------------------------
%11-------------------------------------------------------------------
\vspace*{\fill}
\begin{multicols}{2}
\begin{glyph}{Diagram (図)}{diagram}\label{glyph:diagram}
Diagram/Drawing.\\\hline
By extension comes the idea of ``planning''.
\end{glyph}

\begin{comps}
\comprtwocone & 囗 + 斗 \\\hline
\comprthreecone & enlosure + volume measure  \\\hline
\end{comps}

\begin{remglyph}
    The \hooglig{diagram} depicts an \remcom{enclosure} of \remcom{volume
    measurements}.\\\hline
\end{remglyph}

\begin{prons}
\pronsrtwocone & \textbf{ズ; ト (ZU; TO)} & --- \\\hline
    \pronsrthreecone &---  &はか (haka) \\\hline
\end{prons}

\begin{remprons}
    The \comon{top} \comon{zucchini} \hooglig{diagram} inspired the \ucomkun{haka}.\\\hline
\end{remprons}

%{\middel}
% &  ()\\\hline
\begin{somme}
    図工  & diagram + craft = \textit{manual arts}  & ズ{\middel}コウ (ZU{\middel}K\={O}) \\\hline
    心電図  & heart + electric + diagram = \textit{electrocardiogram}  & シン{\middel}デン{\middel}ズ (SHIN{\middel}DEN{\middel}ZU)  \\\hline
    \notyet{海図}  & sea + diagram = \textit{nautical chart}  & カイ{\middel}ズ (KAI{\middel}ZU)  \\\hline
    \notyet{意図}  & mind + diagram = \textit{(one's) aim}  & イ{\middel}ト (I{\middel}TO) \\\hline
\end{somme}

\begin{sentence}
    \ruby{\underline{心電図}}{シンデンズ}{\;}\underline{モニタアー}\ruby{を}{お}{\;}\ruby{\underline{見}}{み}\underline{せて}{\;}\ruby{\underline{下}}{くだ}\underline{さい}。\\\hline
    SHIN{\middel}DEN{\middel}ZU monit\={a} o mi{\middel}sete kuda{\middel}sai.\\\hline
    (electrocardiogram) (monitor) (show) (please) \\\hline
    =\textit{Please show me the electrocardiogram monitor.}\\\hline
\end{sentence}
\end{multicols}
\vspace*{-\baselineskip}
\vhrulefill{5pt}
%%--------------------------------------------------------------------
%12-------------------------------------------------------------------
\begin{multicols}{2}
\begin{glyph}{Stand (立)}{stand}\label{glyph:stand}
Stand/Rise.
\end{glyph}

\begin{comps}
\comprtwocone & 立 \\\hline
    \comprthreecone & Stand (\ref{glyph:stand}) \\\hline
\end{comps}

\begin{remglyph}
\\\hline
\end{remglyph}

\begin{prons}
\pronsrtwocone & \textbf{リツ (RITSU)} & \textbf{た (ta)}\\\hline
    \pronsrthreecone &リュウ (RY\={U}) & --- \\\hline
    \rye{3}{Consistent with other on-yomi ending in ツ, RITSU ``doubles up''
    with unvoiced consonant sounds that follow it, as the fifth example below
    demonstrates.\\\hline
    It can also absorb its accompanying hiragana, as the final example
    shows.}
\end{prons}

\begin{remprons}
    I cannot \hooglig{stand} \comon{leets}.  It is a \ucomon{leukemia} \comkun{tunneling} its way into our society.\\\hline
\end{remprons}

%{\middel}
% &  ()\\\hline
\begin{somme}
    立つ (intr) & \textit{to stand}  & た{\middel}つ (ta{\middel}tsu) \\\hline
    立てる (tr) & \textit{to set up}  & た{\middel}てる (ta{\middel}teru) \\\hline
    自立  & self + stand = \textit{independence}  & ジ{\middel}リツ (JI{\middel}RITSU)  \\\hline
    中立  & middle + stand = \textit{neutrality}  & チュウ{\middel}リツ (CH\={U}{\middel}RITSU) \\\hline
    立春  & stand + spring = \textit{first day of spring}  &  リッ{\middel}シュン
    (RIS{\middel}SHUN) \\\hline
    立ち入る  & stand + enter = \textit{to go into}  & た{\middel}ち い{\middel}る (ta{\middel}chi i{\middel}ru) \\\hline
    夕立  & evening + stand = \textit{late afternoon (rain) shower}  & ゆう{\middel}だち
    (y\={u}{\middel}dachi)\\\hline
\end{somme}

\begin{sentence}
    \ruby{\underline{女}}{おんな}の\ruby{\underline{人}}{ひと}が\ruby{\underline{美}}{うつく}しい\ruby{\underline{月}}{つき}の\ruby{\underline{下}}{した}に\ruby{\underline{立}}{た}\underline{っています}。\\\hline
    onna no hito ga utsuku{\middel}shii tsuki no shita ni tat{\middel}te imasu.\\\hline
    (woman) (person) (beautiful) (moon) (under) (is standing)\\\hline
    =\textit{The woman is standing under a beautiful moon.}\\\hline
\end{sentence}
\end{multicols}
\vspace*{-\baselineskip}
\vhrulefill{5pt}
\vspace*{\fill}
\pagebreak
%%--------------------------------------------------------------------
%13-------------------------------------------------------------------
\vspace*{\fill}
\begin{multicols}{2}
\begin{glyph}{Parent (親)}{parent}\label{glyph:parent}
``Parent'', together with the senses of intimacy and closeness.
\end{glyph}

\begin{comps}
\comprtwocone & 立 + 木 + 見 \\\hline
    \comprthreecone & stand (\ref{glyph:stand}) + tree (\ref{glyph:tree}) + see
    (\ref{glyph:see})\\\hline
\end{comps}

\begin{remglyph}
    ``\remcom{Stand} on a \remcom{tree}, What do you \remcom{see}?, Your
    \hooglig{parents}, your parents!''\\\hline
\end{remglyph}

\begin{prons}
\pronsrtwocone & \textbf{シン (SHIN)} & \textbf{おや (oya)}\\\hline
    \pronsrthreecone &---  &した (shita) \\\hline
\end{prons}

\begin{remprons}
    The \hooglig{parent} \ucomkun{she tugged} \comon{sheened} with animosity as
    she yelled ``\comkun{Oi ya}! come here''.\\\hline
\end{remprons}

%{\middel}
% &  ()\\\hline
\begin{somme}
    親  & \textit{parent}  & おや (oya) \\\hline
    父親  & father + parent = \textit{father}  & ちち{\middel}おや (chichi{\middel}oya) \\\hline
    親子  & parent + child = \textit{parent and child}  &おや{\middel}こ (oya{\middel}ko) \\\hline
    肉親  & meat + parent = \textit{blood relative}  &ニク{\middel}シン  (NIKU{\middel}SHIN) \\\hline
    \notyet{母親}  & mother + child = \textit{mother}  &はは{\middel}おや (haha{\middel}oya) \\\hline
    \notyet{親切}  & parent + cut = \textit{kindhearted}  &シン{\middel}セツ (SHIN{\middel}SETSU)  \\\hline
\end{somme}

\begin{sentence}
    \ruby{\underline{中川}}{なかがわ}\underline{さん}の\ruby{\underline{父親}}{ちちおや}\ruby{は}{わ}\ruby{\underline{山}}{やま}が\ruby{\underline{大好}}{ダイスキ}\underline{き}です。\\\hline
    Naka{\middel}gawa-san no chichi{\middel}oya wa yama ga DAI{\middel}suki desu.\\\hline
    (Nakagawa-san) (father) (mountain) (really likes)\\\hline
    =\textit{Nakagawa-san's father really likes the mountains.}\\\hline
\end{sentence}
\end{multicols}
\vspace*{-\baselineskip}
\vhrulefill{5pt}
%%--------------------------------------------------------------------
%14-------------------------------------------------------------------
\begin{multicols}{2}
\begin{glyph}{Star (星)}{star}\label{glyph:star}
Star.
\end{glyph}

\begin{comps}
\comprtwocone &日 + 生  \\\hline
    \comprthreecone & Sun (\ref{glyph:sun}) + Life (\ref{glyph:life})  \\\hline
\end{comps}

\begin{remglyph}
    When the \remcom{sun} shines over your \remcom{life}, it is impossible to
    not feel like a \hooglig{star}.\\\hline
\end{remglyph}

\begin{prons}
\pronsrtwocone & \textbf{セイ (SEI)} & \textbf{ほし (hoshi)}\\\hline
    \pronsrthreecone &ショウ (SH\={O})  &---  \\\hline
\end{prons}

\begin{remprons}
    \comon{Sailing} according to the \comkun{horsy} \hooglig{stars} will lead us to the \ucomon{shore}.\\\hline
\end{remprons}

%{\middel}
% &  ()\\\hline
\begin{somme}
    星  & \textit{star}  & ほし (hoshi) \\\hline
    水星  & water + star = \textit{Mercury}  & スイ{\middel}セイ (SUI{\middel}SEI) \\\hline
    金星  & gold + star = \textit{Venus}  & キン{\middel}セイ (KIN{\middel}SEI) \\\hline
    火星  & fire + star = \textit{Mars}  & カ{\middel}セイ (KA{\middel}SEI) \\\hline
    木星  & tree + star = \textit{Jupiter}  & モク{\middel}セイ (MOKU{\middel}SEI) \\\hline
    土星  & earth + star = \textit{Saturn}  & ド{\middel}セイ (DO{\middel}SEI) \\\hline
\end{somme}

\begin{sentence}
    \underline{あの}{\;}\ruby{\underline{明}}{あか}るい\ruby{\underline{星}}{ほし}\ruby{は}{わ}\ruby{\underline{金星}}{キンセイ}だと\ruby{\underline{思}}{おも}\underline{い}ますか。\\\hline
    ano aka{\middel}rui hoshi wa KIN{\middel}SEI da to omo{\middel}imasu ka.\\\hline
    (that) (bright) (star) (Venus) (think)\\\hline
    =\textit{Do you think that bright star is Venus?}\\\hline
\end{sentence}
\end{multicols}
\vspace*{-\baselineskip}
\vhrulefill{5pt}
\vspace*{\fill}
\pagebreak
%%--------------------------------------------------------------------
%15-------------------------------------------------------------------
\vspace*{\fill}
\begin{multicols}{2}
\begin{glyph}{Autumn (秋)}{autumn}\label{glyph:autumn}
Autumn.
\end{glyph}

\begin{comps}
\comprtwocone & 禾 + 火  \\\hline
    \comprthreecone & grain ear + Fire (\ref{glyph:fire})  \\\hline
\end{comps}

\begin{remglyph}
    \hooglig{Autumn} being the time of the year when the stubble of \remcom{wheat} was set on \remcom{fire} after harvesting.\\\hline
\end{remglyph}

\begin{prons}
    \pronsrtwocone & \textbf{シュウ (SH\={U})} & \textbf{あき (aki)}\\\hline
\pronsrthreecone & ---  & --- \\\hline
\end{prons}

\begin{remprons}
    \comkun{Akimbo} \hooglig{autumn} \comon{shoes}.\\\hline
\end{remprons}

%{\middel}
% &  ()\\\hline
\begin{somme}
    秋  & \textit{autumn}  & あき (aki) \\\hline
    立秋  & stand + autumn = \textit{First day of Autumn}  & リッ{\middel}シュウ
    (RIS{\middel}SH\={U}) \\\hline
    春夏秋冬  & spring + summer + autumn + winter = \textit{The four seasons
    year around}  & シュン{\middel}カ{\middel}シュウ{\middel}トウ (SHUN{\middel}KA{\middel}SH\={U}{\middel}T\={O}) \\\hline
    \notyet{秋空}  & autumn + empty = \textit{autumn sky}  & あき{\middel}ぞら (aki{\middel}zora)  \\\hline
    \notyet{秋分}  & autumn + part = \textit{autumnal equinox}  & シュウ{\middel}ブン (SH\={U}{\middel}BUN) \\\hline
\end{somme}

\begin{sentence}
    \ruby{\underline{秋}}{あき}に\ruby{は}{わ}\ruby{\underline{日本}}{ニホン}の\ruby{\underline{山}}{やま}が\ruby{\underline{美}}{うつく}\underline{しい}。\\\hline
    aki ni wa NI{\middel}HON no yama ga utsuku{\middel}shii.\\\hline
    (autumn) (Japan) (mountain) (beautiful)\\\hline
    =\textit{Japan's mountains are beautiful in autumn.}\\\hline
\end{sentence}
\end{multicols}
\vspace*{-\baselineskip}
\vhrulefill{5pt}
%%--------------------------------------------------------------------
%16-------------------------------------------------------------------
\begin{multicols}{2}
\begin{glyph}{West (西)}{west}\label{glyph:west}
West.
\end{glyph}

\begin{comps}
\comprtwocone & 西 \\\hline
    \comprthreecone & West (\ref{glyph:west}) \\\hline
\end{comps}

\begin{remglyph}
\\\hline
\end{remglyph}

\begin{prons}
\pronsrtwocone & \textbf{セイ (SEI)} & \textbf{にし (nishi)}\\\hline
    \pronsrthreecone &サイ (SAI) & --- \\\hline
\end{prons}

\begin{remprons}
    We found our \comkun{niche} of \ucomon{psychedellic} drugs by \comon{sailing} \hooglig{west}.\\\hline
\end{remprons}

%{\middel}
% &  ()\\\hline
\begin{somme}
    西  & \textit{west}  & にし (nishi) \\\hline
    西口  & west + mouth = \textit{western exit}  & にし{\middel}ぐち (nishi{\middel}guchi) \\\hline
    北北西  & north + north + west = \textit{north-northwest}  & ホク{\middel}ホク{\middel}セイ
    (HOKU{\middel}HOKU{\middel}SEI) \\\hline
    南西  & south + west = \textit{southwest}  & ナン{\middel}セイ (NAN{\middel}SEI) \\\hline
    南南西  & south + south + west = \textit{south-southwest}  & ナン{\middel}ナン{\middel}セイ
    (NAN{\middel}NAN{\middel}SEI) \\\hline
    北西  & north + west = \textit{northwest}  & ホク{\middel}セイ (HOKU{\middel}SEI)  \\\hline
\end{somme}

\begin{sentence}
    \ruby{\underline{西口}}{にしぐち}から\ruby{\underline{出}}{で}\underline{て}{\;}\ruby{\underline{父}}{ちち}に\ruby{\underline{会}}{あ}\underline{いました}。\\\hline
    nishi{\middel}guchi kara de{\middel}te chichi ni a{\middel}imashita.\\\hline
    (west exit) (exit) (father) (met)\\\hline
    =\textit{(She) left by the west exit and met her father.}\\\hline
\end{sentence}
\end{multicols}
\vspace*{-\baselineskip}
\vhrulefill{5pt}
\vspace*{\fill}
\pagebreak
%%--------------------------------------------------------------------
%17-------------------------------------------------------------------
\vspace*{\fill}
\begin{multicols}{2}
\begin{glyph}{Sword (刀)}{sword}\label{glyph:sword}
Sword.
\end{glyph}

\begin{comps}
\comprtwocone & 刀 \\\hline
    \comprthreecone & Sword (\ref{glyph:sword}) \\\hline
\end{comps}

\begin{remglyph}
\\\hline
\end{remglyph}

\begin{prons}
    \pronsrtwocone & \textbf{トウ (T\={O})} & \textbf{かたな (katana)}\\\hline
\pronsrthreecone & --- & --- \\\hline
\end{prons}

\begin{remprons}
    He was \comon{tortured} with a \comkun{katana} \hooglig{sword}.\\\hline
\end{remprons}

%{\middel}
% &  ()\\\hline
\begin{somme}
    刀  & \textit{sword}  & かたな (katana) \\\hline
    刀工  & sword + craft = \textit{sword maker}  & トウ{\middel}コウ (T\={O}{\middel}K\={O}) \\\hline
    日本刀  & sun + main + sword = \textit{japanese sword}  & ニ{\middel}ホン{\middel}トウ (NI{\middel}HON{\middel}T\={O}) \\\hline
\end{somme}

\begin{sentence}
    \ruby{\underline{先月}}{センゲツ}{\;}\ruby{\underline{東京}}{トウコウ}で\ruby{\underline{日本刀}}{ニホントウ}\ruby{を}{お}\ruby{\underline{買}}{か}\underline{いました}。\\\hline
    SEN{\middel}GETSU T\={O}{\middel}KY\={O} de NI{\middel}HON{\middel}T\={O} o ka{\middel}imashita.\\\hline
    (last month) (Tokyo) (Japanese Sword) (bought)\\\hline
    =\textit{Last month I bought a Japanese sword in Tokyo.}\\\hline
\end{sentence}
\end{multicols}
\vspace*{-\baselineskip}
\vhrulefill{5pt}
%%--------------------------------------------------------------------
%18-------------------------------------------------------------------
\begin{multicols}{2}
\begin{glyph}{Cut (切)}{cut}\label{glyph:cut}
Cut.\\\hline
When the kunyomi is use as a suffix, it can impart a sense of completely
    finishing something, or acting in a decisive manner.  The final example
    provides an example.
\end{glyph}

\begin{comps}
\comprtwocone &匕 + 刀  \\\hline
    \comprthreecone & spoon + Sword (\ref{glyph:sword}) \\\hline
\end{comps}

\begin{remglyph}
    That \remcom{sword} can \hooglig{cut} a \remcom{spoon}.\\\hline
\end{remglyph}

\begin{prons}
\pronsrtwocone & \textbf{セツ (SETSU)} & \textbf{き (ki)}\\\hline
    \pronsrthreecone &サイ (SAI) & --- \\\hline
\end{prons}

\begin{remprons}
    I \comkun{keep} \hooglig{cutting} \comon{sets} into my \ucomon{psyche}.\\\hline
\end{remprons}

%{\middel}
% &  ()\\\hline
\begin{somme}
    \rye{3}{As the fourth example illustrates, this is another kanji that
    behaves like 買.}
    切れる  & \textit{to break by itself}  & き{\middel}れる (ki{\middel}reru) \\\hline
    切る  & \textit{to cut (something)}  & き{\middel}る (ki{\middel}ru) \\\hline
    切り下げる  & cut + lower = \textit{to devalue}  & き{\middel}り さ{\middel}げる (ki{\middel}ri sa{\middel}geru) \\\hline
    切下げ  & cut + lower = \textit{devaluation}  & き{\middel}り さ{\middel}げ (ki{\middel}ri sa{\middel}ge) \\\hline
    親切  & parent + cut = \textit{kindhearted}  & シン{\middel}セツ (SHIN{\middel}SETSU) \\\hline
    大切  & large + cut = \textit{important}  & タイ{\middel}セツ (TAI{\middel}SETSU) \\\hline
    思い切る  & think + cut = \textit{to make up one's mind}  & おも{\middel}い き{\middel}る (omo{\middel}i ki{\middel}ru) \\\hline
\end{somme}

\begin{irrr}
    \rye{3}{As the kunyomi of 切 ``doubles up'' with the kanji following it in
    this next pair of words, these compounds are best learned as irregular
    readings.}
    切手  & cut + hand = \textit{postage stamp}  & きっ{\middel}て (kit{\middel}te) \\\hline
    切符{\notcov}  & cut + symbol = \textit{ticket}  & きっ{\middel}プ (kip{\middel}pu)  \\\hline
\end{irrr}

\begin{sentence}
    \ruby{\underline{水力}}{スイリョク}と\ruby{\underline{火力}}{カリョク}\ruby{は}{わ}\ruby{\underline{大切}}{タイセツ}ですね。\\\hline
    SUI{\middel}RYOKU to KA{\middel}RYOKU wa TAI{\middel}SETSU desu ne.\\\hline
    (hydro power) (thermal power) (important)\\\hline
    =\textit{Hydro and thermal power are important, aren't they?}\\\hline
\end{sentence}
\end{multicols}
\vspace*{-\baselineskip}
\vhrulefill{5pt}
\vspace*{\fill}
\pagebreak
%%--------------------------------------------------------------------
%19-------------------------------------------------------------------
\vspace*{\fill}
\begin{multicols}{2}
\begin{glyph}{Day of the week (曜)}{day_of_the_week}\label{glyph:day_of_the_week}
Day of the week.
\end{glyph}

\begin{comps}
\comprtwocone &日 + 彐 + 彐 + 隹  \\\hline
    \comprthreecone &Sun (\ref{glyph:sun}) + pig heads + small bird  \\\hline
\end{comps}

\begin{remglyph}
    Every \hooglig{day of the week}, the \remcom{small birds} sit on the \remcom{pig heads}.\\\hline
\end{remglyph}

\begin{prons}
    \pronsrtwocone & \textbf{ヨウ (Y\={O})} & ---\\\hline
\pronsrthreecone & --- & --- \\\hline
\end{prons}

\begin{remprons}
    We \comon{yawn} at least once \hooglig{everday of the week}.\\\hline
\end{remprons}

%{\middel}
% &  ()\\\hline
\begin{somme}
    日曜日  & sun + day of the week + sun (day) = \textit{Sunday}  & ニチ{\middel}ヨウ{\middel}び
    (NICHI{\middel}Y\={O}{\middel}bi)\\\hline
    月曜日  & mooon + day of the week + sun (day) = \textit{Monday}  &
    ゲツ{\middel}ヨウ{\middel}び (GETSU{\middel}Y\={O}{\middel}bi) \\\hline
    火曜日  & fire + day of the week + sun (day) = \textit{Tuesday}  & カ{\middel}ヨウ{\middel}び
    (KA{\middel}Y\={O}{\middel}bi)\\\hline
    水曜日  & water + day of the week + sun (day) = \textit{Wednesday}  &
    スイ{\middel}ヨウ{\middel}び (SUI{\middel}Y\={O}{\middel}bi) \\\hline
    木曜日  & tree + day of the week + sun (day) = \textit{Thursday}  &
    モク{\middel}ヨウ{\middel}び (MOKU{\middel}Y\={O}{\middel}bi) \\\hline
    金曜日  & gold + day of the week + sun (day) = \textit{Friday}  &
    キニ{\middel}ヨウ{\middel}び
    (KIN{\middel}Y\={O}{\middel}bi)\\\hline
    土曜日  & earth + day of the week + sun (day) = \textit{Saturday}  &
    ド{\middel}ヨウ{\middel}び (DO{\middel}Y\={O}{\middel}bi) \\\hline
\end{somme}

\begin{sentence}
    \ruby{\underline{土曜日}}{ドヨウび}に\underline{パーティー}が\underline{あります}。\\\hline
    DO{\middel}Y\={O}{\middel}bi ni p\={a}{\middel}t\={i} ga arimasu.\\\hline
    (Saturday) (party) (is)\\\hline
    =\textit{There's a party on Saturday}\\\hline
\end{sentence}
\end{multicols}
\vspace*{-\baselineskip}
\vhrulefill{5pt}
%%--------------------------------------------------------------------
%20-------------------------------------------------------------------
\begin{multicols}{2}
\begin{glyph}{Around (周)}{around}\label{glyph:around}
Around, as in the sense of a lap or circuit.
\end{glyph}

\begin{comps}
\comprtwocone & 冂 +  土 + 口 \\\hline
    \comprthreecone & wilderness + Earth (\ref{glyph:earth}) + Mouth (\ref{glyph:mouth}) \\\hline
\end{comps}

\begin{remglyph}
    \hooglig{Around} the \remcom{earthen} \remcom{wilderness}, your
    \remcom{mouths} cannot be satisfied.\\\hline
\end{remglyph}

\begin{prons}
    \pronsrtwocone & \textbf{シュウ (SH\={U})} & --- \\\hline
\pronsrthreecone & --- & --- \\\hline
\end{prons}

\begin{remprons}
    \comon{Shoe} \hooglig{lap}.\\\hline
\end{remprons}

%{\middel}
% &  ()\\\hline
\begin{somme}
    円周  & circle + around = \textit{circumference}  & エン{\middel}シュウ (EN{\middel}SH\={U}) \\\hline
    周回  & around + rotate = \textit{going around}  & シュウ{\middel}カイ (SH\={U}{\middel}KAI) \\\hline
    一周  & once + around = \textit{once around}  & イッ{\middel}シュウ (IS{\middel}SH\={U}) \\\hline
\end{somme}

\begin{sentence}
    \ruby{\underline{寺田}}{てらだ}\underline{さん}\ruby{は}{わ}\underline{ニュウジランド}\ruby{を}{お}\ruby{\underline{一周}}{イッシュウ}{\;}\underline{しました}。\\\hline
    Tera{\middel}da-san wa Ny\={u}j\={i}rando o IS{\middel}SH\={U} shimashita.\\\hline
    (Terada-san) (New Zealand) (once around) (did)\\\hline
    =\textit{Terada-san did a circuit of New Zealand.}\\\hline
\end{sentence}
\end{multicols}
\vspace*{-\baselineskip}
\vhrulefill{5pt}
%%--------------------------------------------------------------------
%%--------------------------------------------------------------------
\vspace*{\fill}
\pagebreak
